\section*{Contoh}
\addcontentsline{toc}{section}{Contoh}

\begin{contoh}[Aliran token program Subset-C]
Diberikan program
\begin{lstlisting}[language=C]
int main() {
    int a = 5;
    print(a);
    return 0;
}
\end{lstlisting}
Aliran token yang dihasilkan secara konseptual adalah:
\begin{center}
\small
\code{INT} \code{ID(main)} \code{LPAREN} \code{RPAREN} \code{LBRACE}
\code{INT} \code{ID(a)} \code{ASSIGN} \code{INT\_LIT(5)} \code{SEMI}
\code{PRINT} \code{LPAREN} \code{ID(a)} \code{RPAREN} \code{SEMI}
\code{RETURN} \code{INT\_LIT(0)} \code{SEMI} \code{RBRACE} \code{EOF}.
\end{center}
Setiap token membawa lokasi (baris, kolom) untuk keperluan diagnostik.
\end{contoh}

\begin{contoh}[Compiler versus interpreter]
Program yang sama dijalankan pada compiler dan interpreter. Pada
compiler, program sumber dikompilasi menjadi executable, lalu
executable dijalankan berulang kali. Pada interpreter, program sumber
dibaca dan dieksekusi setiap kali dijalankan. Jika program dijalankan
100 kali, compiler membayar biaya kompilasi sekali, sedangkan
interpreter membayar biaya parsing dan evaluasi 100 kali. Namun, untuk
program yang hanya dijalankan sekali, interpreter sering lebih cepat
karena tidak perlu menghasilkan kode target.
\end{contoh}

\begin{contoh}[Activation record]
Misalkan fungsi \code{faktorial(3)} dipanggil. Frame pertama memuat
\code{n = 3} dan alamat kembali ke pemanggil. Karena \code{n > 1},
fungsi memanggil \code{faktorial(2)}, mendorong frame kedua dengan
\code{n = 2}. Proses berlanjut hingga \code{faktorial(1)}. Setelah
basis tercapai, nilai 1 dikembalikan, lalu frame dilepas satu per
satu sambil mengalikan hasilnya: \(1 \times 2 \times 3 = 6\).
\end{contoh}

\begin{contoh}[Memilih strategi kompilasi]
Sebuah tim mengembangkan bahasa domain khusus untuk perangkat
mikrokontroler dengan memori sangat terbatas. Mereka memilih
kompilator satu lintasan yang langsung menghasilkan kode mesin.
Keputusan ini masuk akal karena memori kecil dan optimasi agresif
tidak menjadi prioritas. Sebaliknya, tim yang mengembangkan bahasa
umum untuk banyak platform memilih front-end bersama dan beberapa
back-end, sehingga investasi pada analisis dan IR terbayar.
\end{contoh}
